\exercisesection{完全交}

\begin{enumerate}
\def\labelenumi{\arabic{enumi})}
\item
  设 \(k\) 为域，A 为多项式环 \(k[X,Y,Z]\) 对由 \(Z^2,ZX,Z(Y-1)\) 生成的理想的商环，\(x\) 与 \(y\) 为 X 与 Y 在 \(\Lambda\) 中的类。证明序列 \((x,y)\) 对 A 是完全割的但不是 A-正则的，而序列 \((y,x)\) 却是 A-正则的。
\item
  设 A 为正则局部环，I 为 A 的一个理想。
\end{enumerate}

\begin{enumerate}
\def\labelenumi{\alph{enumi})}
\tightlist
\item
  若 ht(I) = 1，证明下列条件等价：
\end{enumerate}

\begin{enumerate}
\def\labelenumi{(\roman{enumi})}
\item
  理想 I 是主理想；
\item
  A/I 是 Gorenstein 环；
\item
  A/I 是 Macaulay 环。
\end{enumerate}

\begin{enumerate}
\def\labelenumi{\alph{enumi})}
\setcounter{enumi}{1}
\item
  假设 \(ht(I)=2\)。为使 A/I 是 Gorenstein 环，必须且只需理想 I 是完全割的（设 \((L,p)\) 为 I 的一个极小投射分解；注意 \(L_{1}\) 的秩为 1，从而 \(L_{0}\) 的秩为 2）。
\item
  设 \(k\) 为交换域；取 \(A = k[[X,Y,Z]]\) 与 \(I = (XY,YZ,ZX, X^2 - Y^2, Y^2 - Z^2)\)。证明 \(I\) 是一个高度为 3 的不完全割的理想，而 \(A/I\) 是 Gorenstein 环。
\end{enumerate}

\begin{enumerate}
\def\labelenumi{\arabic{enumi})}
\setcounter{enumi}{2}
\tightlist
\item
  设 R 为正则局部环，I 为 R 的理想，A 为环 R/I。设 \(\mathbf{x} = (x_{1}, \ldots, x_{d})\) 为 R 的一个坐标系，\(\mathbf{y} = (y_{1}, \ldots, y_{m})\) 为 I 的一个极小生成系，于是 \(m = [I/\mathfrak{m}_{R}\mathbf{l} : \kappa_{R}]\)。
\end{enumerate}

\begin{enumerate}
\def\labelenumi{\alph{enumi})}
\item
  当 I 包含于 \(\mathfrak{m}_{\mathrm{R}}^{2}\) 时，定义 \(\mathrm{H}_1(\mathbf{x},\mathrm{A})\) 到 \(\mathrm{I} / \mathfrak{m}_{\mathrm{R}}\mathrm{I}\) 上的一个典范同构（考虑 Koszul 复形的正合列）。
\item
  为使 \(x_{i}\) 在 A 中的像构成 \(m_{A}\) 的极小生成系，必须且只需 I 包含于 \(m_{R}^{2}\)。
\item
  证明 \(\kappa_{\mathrm{A}}\)-向量空间 \(H_1(\mathbf{x}, A)\) 的维数等于 \(m - d + [\mathfrak{m}_{\mathrm{A}} / \mathfrak{m}_{\mathrm{A}}^{2} : \kappa_{\mathrm{A}}]\)。
\end{enumerate}

4) 设 A 为 Noether 局部环。令 \(\delta(A) = [\mathfrak{m}_{A}/\mathfrak{m}_{A}^{2}:\kappa_{A}] - \dim(A)\)（参见 § 4, 第 5 节）。

\begin{enumerate}
\def\labelenumi{\alph{enumi})}
\item
  证明 \(\kappa_{\Lambda}\)-向量空间 \(H_{i}(\mathbf{z}, A)\)（其中 \(\mathbf{z}\) 为 \(\mathfrak{m}_{\mathrm{A}}\) 的一个极小生成系）的维数与 \(\mathbf{z}\) 的选取无关；记之为 \(h_{i}(\mathrm{A})\)。
\item
  证明等式 \(h_i(\mathrm{A}) = h_i(\widehat{\mathrm{A}})\) 对每个 \(i\) 成立。
\item
  证明有 \(h_1(A) \geqslant \delta(A)\)，且等号成立当且仅当 \(A\) 是完全交环（第 2 节的注 2；化为 \(A\) 完备的情形，并把 \(A\) 表为商 \(R/I\)（其中 \(R\) 为正则局部环且 \(I \subset m_R^2\)）后应用习题 4）。
\item
  设 \((x_{1},\ldots,x_{p})\) 为 \(m_{A}\) 中元素的生成理想 J 的一个完全割序列。证明等式 \(h_{1}(A/J)-\delta(A/J)=h_{1}(A)-\delta(A)\)（化为 p=1 的情形，并按 \(x_{1}\) 属于或不属于 \(m_{A}^{2}\) 分两种情形讨论）。
\end{enumerate}

\begin{enumerate}
\def\labelenumi{\arabic{enumi})}
\setcounter{enumi}{4}
\tightlist
\item
  设 R 为正则局部环，I 为 R 的理想，A 为环 R/I。为使 A 是完全交环，必须且只需理想 I 是完全割的（利用习题 4 c) 与 3 c)）。
\end{enumerate}

¶ 6) 设 A 为 Noether 局部环，I 为 A 的理想。

\begin{enumerate}
\def\labelenumi{\alph{enumi})}
\tightlist
\item
  证明下列条件等价：
\end{enumerate}

\begin{enumerate}
\def\labelenumi{(\roman{enumi})}
\item
  理想 I 是完全割的；
\item
  A/I-模 I/I² 是自由的且 dp\_A(A/I) \textless{} +∞。
\end{enumerate}

在条件 (ii) 下，由 § 3 习题 2 推出存在 I 中的可消去元 \(x\)，使 \(x \notin \mathfrak{m}_{\mathrm{A}}\mathrm{I}\)；对 \(\mathrm{I} / \mathrm{I}^2\) 的维数作归纳论证。）

\begin{enumerate}
\def\labelenumi{\alph{enumi})}
\setcounter{enumi}{1}
\item
  证明 A 中使理想 \(I_{p}\) 为完全割的素理想 p 所成的集合在 Spec(A) 中是开的。
\item
  设 B 为可表现环；证明 B 中使 \(B_{q}\) 为完全交环的素理想 q 所成的集合在 Spec(B) 中是开的。
\end{enumerate}

\begin{enumerate}
\def\labelenumi{\arabic{enumi})}
\setcounter{enumi}{6}
\tightlist
\item
  设 A 为环，B 为一个 A-代数；假设 A-模 B 是有限型自由的。在以下习题中，用 B* 表示 A-模 Hom \(_{A}\)(B, A)，并赋予其自然的 B-模结构。
\end{enumerate}

\begin{enumerate}
\def\labelenumi{\alph{enumi})}
\item
  设 \((e_i)_{i\in I}\) 为 A-模 B 的一个基，\((e_i^*)_{i\in I}\) 为 \(\mathbf{B}^*\) 的对偶基。证明 \(\mathrm{Tr}_{\mathrm{B / A}}\)（\(\mathbf{B}^*\) 的元素，A, III, 第~110 页）等于 \(\sum_{j}e_{i}e_{i}^{*}\)。
\item
  假设 B-模 \(B^{*}\) 是秩 1 自由的；设 \(\ell\) 为构成该模的一个基的元素，\(\delta\) 为 B 中使 \(Tr_{B/A} = \delta\ell\) 的元素。B 的理想 \(\delta B\) 不依赖于 \(\ell\) 的选取；称之为 B 在 A 上的微差理想。证明 \(N_{B/A}(\delta)\) 是 B 在 A 上的判别式理想的一个生成元（A, III, 第 115 页；计算 B 的基（\(e_{i}\)）在 \(\Lambda\) 上的判别式，利用 \(\delta\) 在此基下的乘法矩阵）。
\end{enumerate}

\begin{enumerate}
\def\labelenumi{\arabic{enumi})}
\setcounter{enumi}{7}
\tightlist
\item
  设 A 为环，P 为次数 \(n\) 的 A{[}X{]} 首一多项式，B 为 A-代数 A{[}X{]}/(P)。用 \(x\) 表示 X 在 B 中的类。设 \(\ell: \mathrm{B} \to \mathrm{A}\) 为满足 \(\ell(x^{n-1}) = 1\)、且 \(\ell(x^i) = 0\)（\(0 \leqslant i \leqslant n-2\)）的 A-线性形式；用 \(\tilde{\ell}: \mathrm{B}[\mathrm{X}] \to \mathrm{A}[\mathrm{X}]\) 表示 A-映射{[}X{]}-线性映射（由 \(\ell\) 经纯量扩张导出）。
\end{enumerate}

\begin{enumerate}
\def\labelenumi{\alph{enumi})}
\item
  设 \(\sigma : B[X] \to A[X]\) 为 \(A[X]\)-线性映射，满足 \(\sigma(x^i) = X^i\)（\(i = 0, \ldots, n - 1\)）。证明公式 \(P\tilde{\ell}(Q) = \sigma((X - x)Q)\) 对每个元素 \(Q\)（属于 \(B[X]\)）成立（对 \(Q = 1, \ldots, x^{n-1}\) 验证之）。
\item
  设 \(\mathrm{P}^{*}(\mathrm{X}) = b_{n-1}\mathrm{X}^{n-1} + \ldots + b_{0}\) 为 B{[}X{]} 中使 \(\mathrm{P}(\mathrm{X}) = (\mathrm{X} - x)\mathrm{P}^{*}(\mathrm{X})\) 在 B{[}X{]} 中成立的多项式。证明 \((b_{0}\ell, \ldots, b_{n-1}\ell)\) 是 \(B^{*}\) 的与 B 的基 \((1, \ldots, x^{n-1})\) 对偶的基（对多项式 \(x^{i}\mathrm{P}^{*}\) 应用 a)）。由此推出 B-模 \(B^{*}\) 是秩 1 自由的，且 \(\ell\) 构成它的一个基。
\item
  证明等式 \(\mathrm{Tr}_{\mathrm{B / A}} = \mathrm{P}'(x)\ell\)（在 \(\mathrm{B}^*\) 中；参见习题 7 a)）。
\end{enumerate}

\begin{enumerate}
\def\labelenumi{\arabic{enumi})}
\setcounter{enumi}{8}
\tightlist
\item
  设 A 为环，n 为整数，S 为代数 \(A[X_{1},\ldots,X_{n}]\)（相应地，\(A[[X_{1},\ldots,X_{n}]]\)），a 为 S 的一个由对 S 完全割的序列 \(\mathbf{P}=(\mathbf{P}_{1},\ldots,\mathbf{P}_{n})\) 生成的理想。令 \(B=S/a\)，并用 \(\pi\) 表示 S 到 S/a 上的典范同态。假设 A-模 B 是有限型自由的。我们旨在证明 B-模 \(B^{*}\) 是秩 1 自由的，且 B 在 A 上的微差理想由 \(\pi(\det(\frac{\partial P_{i}}{\partial X_{j}}))\) 生成。
\end{enumerate}

\begin{enumerate}
\def\labelenumi{\alph{enumi})}
\item
  令 \(C = S \otimes_{A} B\)，并令 \(\xi_{i} = X_{i} \otimes 1 - 1 \otimes \pi(X_{i})\)（\(1 \leqslant i \leqslant n\)）。用 \(\varphi : C \to B\) 表示满足 \(\varphi(P \otimes b) = \pi(P)b\) 的 A-代数同态。证明 \(\varphi\) 的核由序列 \(\boldsymbol{\xi} = (\xi_{1}, \ldots, \xi_{n})\) 生成，且该序列对 C 是完全割的。
\item
  设 \((c_{ij}) \in \mathbf{M}_{n}(\mathbf{C})\) 为满足 \(P_{i} \otimes 1 = \sum_{j} c_{ij} \xi_{j}\) 的矩阵。证明等式 \(\pi(\frac{\partial P_{i}}{\partial X_{j}}) = \varphi(c_{ij})\)。
\item
  矩阵 \((c_{ij})\) 定义复形的一个态射 \(\mathbf{K}_{\bullet}(\mathbf{P},\mathrm{C})\to\mathbf{K}_{\bullet}(\boldsymbol{\xi},\mathrm{C})\)（A, X, 第~151 页），从而给出 S-模复形的一个态射 \(u:\mathbf{K}_{\bullet}(\mathbf{P},\mathrm{S})\to\mathbf{K}_{\bullet}(\boldsymbol{\xi},\mathrm{C})\)。同态 \(u_{n}:S\to C\) 把 \(1_{S}\) 映到 C 的元素 \(d=\det(c_{ij})\)。
\item
  设 \(v: \mathrm{Hom}_{\mathrm{S}}(\mathrm{C}, \mathrm{S}) \otimes_{\mathrm{C}} \mathrm{B} \to \mathrm{B}\) 为满足 \(v(f \otimes 1) = \pi(f(d))\) 的 B-线性同态；证明 \(v\) 是同构（把 \(v\) 与 \(\mathrm{H}^n(\mathrm{Homgr}_{\mathrm{S}}(u, 1_{\mathrm{S}}))\) 等同，并注意复形 \(\mathbf{K}_{\bullet}(\mathbf{P}, \mathrm{S})\) 与 \(\mathbf{K}_{\bullet}(\boldsymbol{\xi}, \mathrm{C})\) 中的每一个都给出 B 的一个由有限型自由 S-模组成的分解）。

\item
  考虑复合映射 \(t: \mathrm{Hom}_{\mathrm{A}}(\mathrm{B}, \mathrm{A}) \to \mathrm{Hom}_{\mathrm{S}}(\mathrm{C}, \mathrm{S}) \to \mathrm{B}\)，其中第一个箭头由纯量扩张得到，第二个由 \(v\) 得到。证明 \(t\) 是 B-线性且双射的（把 \(\mathrm{Hom}_{\mathrm{S}}(\mathrm{C}, \mathrm{S})\) 与 \(\mathrm{Hom}_{\mathrm{A}}(\mathrm{B}, \mathrm{A}) \otimes_{\mathrm{B}} \mathrm{C}\) 等同）。
\item
  证明公式 \(t(\mathrm{Tr}_{\mathrm{B / A}}) = \varphi (d)\)（用 \(d\) 与 \(\mathrm{Tr}_{\mathrm{B / A}}\) 由 B 的一个基及 \(\mathrm{B}^*\) 的对偶基表出，参见习题 7, a)）。
\item
  由 b)、e) 与 f) 推出 B 在 A 上的微差理想由 \(\pi(\det(\frac{\partial P_i}{\partial X_j}))\) 生成。
\end{enumerate}

\begin{enumerate}
\def\labelenumi{\arabic{enumi})}
\setcounter{enumi}{9}
\tightlist
\item
  设 A 为离散赋值环，v 为其赋范赋值，K 为其分式域。考虑一个 Noether A-代数 B，它带有 A-代数同态 \(\varepsilon:B\to A\)。用 I 表示 \(\varepsilon\) 的核，J 表示其在 B 中的零化子。
\end{enumerate}

\begin{enumerate}
\def\labelenumi{\alph{enumi})}
\tightlist
\item
  证明下列条件等价：
\end{enumerate}

\begin{enumerate}
\def\labelenumi{(\roman{enumi})}
\item
  \(\Lambda\)-模 \(1 / I^2\) 是有限长度的；
\item
  \(K \otimes_{A} B\) 在极大理想 \(K \otimes_{A} I\) 处的局部环（作为 K-代数）同构于 K；
\item
  我们有 \(\mathrm{K}\otimes_{\mathrm{A}}\mathrm{B} = (\mathrm{K}\otimes_{\mathrm{A}}\mathrm{I})\oplus (\mathrm{K}\otimes_{\mathrm{A}}\mathrm{J})\)
\end{enumerate}

以下假设这些条件满足；用 \(\gamma(B, \varepsilon)\)（或简作 \(\gamma(B)\)）表示 A-模 \(I/I^{2}\) 的长度，用 \(\eta(B, \varepsilon)\)（或简作 \(\eta(B)\)）表示 A-模 \(A/\varepsilon(J)\) 的长度。

\begin{enumerate}
\def\labelenumi{\alph{enumi})}
\setcounter{enumi}{1}
\item
  设 b 为 B 的一个含于 I 的理想；考虑 A-代数 \(B' = B/b\) 以及同态 \(\varepsilon': B' \to A\)（由 \(\varepsilon\) 导出）。证明 \(B'\) 也满足 a) 的条件，且有 \(\gamma(B', \varepsilon') \leqslant \gamma(B, \varepsilon)\) 与 \(\eta(B', \varepsilon') \leqslant \eta(B, \varepsilon)\)。
\item
  假设 A-模 B 是有限型自由的。证明 J 是秩 1 自由的 A-模，是 B 的一个直和因子，且 \(\mathrm{Tr}_{\mathrm{B / A}}(x) = \varepsilon (x)\) 对每个 \(x\in J\) 成立。
\item
  进而假设 B 是 Gorenstein 环；用 \(\ell\) 表示 B-模 B* 的一个基（§ 3, 习题 5），用 δ 表示 B 在 A 上的微差理想的一个生成元（习题 7, b)）。证明等式 η(B) = v(ε(δ))（证明 \(\ell\)(J) 等于 A，并利用 c)）。
\item
  置于 \(b\) 的情形，且 \(\mathfrak{b} \neq 0\)；假设 A-代数 B 与 B' 都是有限且平坦的，且 B 是 Gorenstein 环。证明有 \(\eta(B') < \eta(B)\)（注意 J 在 B/ \(\mathfrak{m}_{A}\) B 中的像是 B/ \(\mathfrak{m}_{A}\) B 的底座，故含于 \(\mathfrak{b}/\mathfrak{m}_{A}\mathfrak{b}\)）。
\end{enumerate}

¶ 11) 保留上一习题的记号；假设环 A 与 B 都是局部的且完备的。

\begin{enumerate}
\def\labelenumi{\alph{enumi})}
\item
  证明 B 同构于形式幂级数代数 \(S = A[[X_{1}, \ldots, X_{n}]]\) 对一个含于 \((X_{1}, \ldots, X_{n})\) 的理想 a 的商；可取 \(n = \dim_{\kappa_{A}}(\mathfrak{n}/\mathfrak{n}^{2})\)，其中 n 为 \(m_{B}\) 在 \(B/m_{A}B\) 中的像。
\item
  证明下列条件等价：
\end{enumerate}

\begin{enumerate}
\def\labelenumi{(\roman{enumi})}
\item
  B 是有限平坦 A-代数且是完全交环；
\item
  理想 a 由一个序列 \((\mathrm{P}_{1},\ldots,\mathrm{P}_{n})\) 生成，该序列对 \(\kappa_{A}\otimes_{A}S\) 完全割。
\end{enumerate}

（参见习题 5。）

\begin{enumerate}
\def\labelenumi{\alph{enumi})}
\setcounter{enumi}{2}
\item
  若这些条件满足，证明有 \(\eta(B) = \gamma(B) = v(\varepsilon(\det(\frac{\partial P_i}{\partial X_j})))\)（利用习题 9 与 10，\(d\)）。
\item
  证明可以找到 a 中元素的一个序列 \((f_{1},\ldots,f_{n})\)，它对 \(\kappa_{A}\otimes_{A}S\) 完全割，使 A-代数 \(\mathrm{B}^{\prime}=\mathrm{S}/(f_{1},\ldots,f_{n})\) 满足 \(\gamma(\mathrm{B}^{\prime})=\gamma(\mathrm{B})\)。
\end{enumerate}

（由正合列 \(\mathfrak{a}/\mathfrak{a}^{2}\to\Omega_{\mathrm{A}}(\mathrm{S})\otimes_{\mathrm{S}}\mathrm{B}\to\Omega_{\mathrm{A}}(\mathrm{B})\to0\)（A, III, 第~137 页）导出 A-模的正合列 \(\mathfrak{a}/\mathfrak{a}^{2}\otimes_{\mathrm{S}}\mathrm{A}\to\Omega_{\mathrm{A}}(\mathrm{S})\otimes_{\mathrm{S}}\mathrm{A}\to\mathrm{I}/\mathrm{I}^{2}\to0\)；取元素 \(f_{1},\ldots,f_{n}\)（取自 \(\mathfrak{a}\)），使其在 \(\Omega_{\mathrm{A}}(\mathrm{S})\otimes_{\mathrm{S}}\Lambda\) 中的像构成 \(\mathfrak{a}/\mathfrak{a}^{2}\otimes_{\mathrm{S}}\mathrm{A}\) 的像在 A 上的一个基。）

\begin{enumerate}
\def\labelenumi{\alph{enumi})}
\setcounter{enumi}{4}
\tightlist
\item
  于是，利用习题 10, e)，我们有 \(\gamma(B) \geqslant \eta(B)\)，且等号成立当且仅当 B 是完全交环。
\end{enumerate}

